\refstepcounter{chapter}
\chapter*{\uppercase{APPENDIX}}
\hypertarget{appendix-heading}{}
\makeatletter
\def\@currentHref{appendix-heading}
\addcontentsline{toc}{chapter}{APPENDIX}
\makeatother
\vspace{-0.4cm}
\hypertarget{appendix-a-heading}{}
\begin{center}
{\Large\bfseries APPENDIX \thechapter\\[0.25em]
Agent Catalogue and Internal Workflows}
\end{center}
\makeatletter
\def\@currentHref{appendix-a-heading}
\addcontentsline{toc}{chapter}{\protect\numberline{\thechapter}\uppercase{Agent Catalogue and Internal Workflows}}
\makeatother
\label{app:agents}

This appendix documents the agent implementations in
\texttt{server/kernel/Agents/agents/}. The diagrams describe the internal
responsibility of each agent and its actual tool or model interaction. The
planner, scheduler, global execution state, and frontend are intentionally not
shown here because they are documented in the system-design chapter.
The registered agents and their implemented responsibilities are summarised in
Table~\ref{tab:agent-catalogue}.

\begin{figure}[!htbp]
\centering
\begin{tabularx}{\textwidth}{>{\raggedright\arraybackslash}p{0.24\textwidth}>{\raggedright\arraybackslash}X}
\toprule
Agent & Implemented responsibility \\
\midrule
A1 & System information and resource inspection. \\
A2 & Date, time, and timezone queries. \\
A3 & Factual and academic research with sources. \\
A4 & File and document operations with deletion confirmation. \\
Content Creator & Text drafting, rewriting, and formatting. \\
Email & Email dispatch and delivery. \\
Weather & Current weather and forecast retrieval. \\
Loan EMI & EMI and amortisation calculation. \\
Crypto PnL & Cryptocurrency trade PnL and ROI reporting. \\
Code Runner & Python generation, execution, capture, and repair. \\
Agent Creator & Creation and registration of new specialised agents. \\
\bottomrule
\end{tabularx}
\caption{AgentOS agent catalogue}
\label{tab:agent-catalogue}
\end{figure}

\newcommand{\AgentArchitecture}[6]{%
\begin{figure}[!htbp]
\centering
\resizebox{\textwidth}{!}{%
\begin{tikzpicture}[
  node distance=11mm and 13mm,
  input/.style={draw, rounded corners, align=center, minimum height=15mm, text width=36mm, fill=blue!8},
  core/.style={draw, rounded corners, align=center, minimum height=20mm, text width=46mm, fill=orange!22, very thick},
  tool/.style={draw, rounded corners, align=center, minimum height=20mm, text width=46mm, fill=green!12},
  output/.style={draw, rounded corners, align=center, minimum height=15mm, text width=36mm, fill=blue!8},
  line/.style={-{Latex[length=2mm]}, thick}
]
\node[input] (request) {Task and\\previous context\\\scriptsize #2};
\node[core, right=of request] (controller) {#1\\\scriptsize agent-specific instructions};
\node[tool, right=of controller] (tools) {#4\\\scriptsize tool calls or model generation};
\node[output, right=of tools] (result) {Validated result\\\scriptsize #5};
\draw[line] (request) -- (controller);
\draw[line] (controller) -- (tools);
\draw[line] (tools) -- (result);
\end{tikzpicture}}
\caption{Internal architecture of #1}
\label{fig:#6-architecture}
\end{figure}%
}

\newcommand{\AgentFlow}[5]{%
\begin{figure}[!htbp]
\centering
\resizebox{0.92\textwidth}{!}{%
\begin{tikzpicture}[
  node distance=9mm,
  inputflow/.style={draw, rounded corners, align=center, minimum height=16mm, text width=45mm, fill=blue!10, very thick},
  controlflow/.style={draw, rounded corners, align=center, minimum height=16mm, text width=45mm, fill=gray!12},
  capabilityflow/.style={draw, rounded corners, align=center, minimum height=16mm, text width=45mm, fill=orange!18},
  outputflow/.style={draw, rounded corners, align=center, minimum height=16mm, text width=45mm, fill=green!12, very thick},
  line/.style={-{Latex[length=2mm]}, thick},
  font=\small
]
\node[inputflow] (one) {1. Input Specification\\#2};
\node[controlflow, right=of one] (two) {2. Agent Control Logic\\construct instructions and validate context};
\node[capabilityflow, right=of two] (three) {3. Capability Execution\\#3};
\node[outputflow, right=of three] (four) {4. Result Synthesis\\#4};
\draw[line] (one) -- (two);
\draw[line] (two) -- (three);
\draw[line] (three) -- (four);
\end{tikzpicture}}
\caption{Internal flow of #1}
\label{fig:#5-flow}
\end{figure}%
}

% The appendix retains one academically styled execution-flow diagram per
% agent. The earlier generic architecture diagrams are intentionally omitted.
\renewcommand{\AgentArchitecture}[6]{}
\renewcommand{\AgentFlow}[5]{}

\tikzset{
  logicstep/.style={draw, rounded corners, align=center, minimum height=13mm, text width=38mm, fill=blue!8},
  logicdata/.style={draw, rounded corners, align=center, minimum height=13mm, text width=38mm, fill=green!12},
  logicdecision/.style={draw, diamond, aspect=2, align=center, inner sep=2pt, fill=yellow!20},
  logicwarning/.style={draw, rounded corners, align=center, minimum height=13mm, text width=38mm, fill=red!10},
  logicline/.style={-{Latex[length=2mm]}, thick},
  logiclabel/.style={font=\small},
  every picture/.style={scale=0.84, transform shape}
}

\clearpage
\section{Agent A1: System Information Agent}

The A1 implementation creates a LangChain agent with the single registered
tool \texttt{get\_system\_info}. The tool reports the operating-system name
and release, machine architecture, and processor descriptor using Python's
\texttt{platform} module. The agent therefore provides a compact host
identity report; process enumeration and detailed resource monitoring are not
implemented by the current tool.
The system-information processing logic is shown in Figure~\ref{fig:a1-logic}.
\AgentArchitecture{Agent A1}{system-state question}{system inspection and analysis}{\texttt{A1\_TOOLS}}{clear system report}{a1}
\begin{figure}[!htbp]\centering
\begin{tikzpicture}[node distance=8mm, logiclabel]
\node[logicstep] (request) {System information request};
\node[logicstep, below=of request] (call) {Invoke \texttt{get\_system\_info}};
\node[logicdata, below left=of call] (os) {Operating system and release};
\node[logicdata, below=of call] (arch) {Machine architecture};
\node[logicdata, below right=of call] (cpu) {Processor descriptor};
\node[logicstep, below=18mm of arch] (report) {Assemble system information report};
\draw[logicline] (request) -- (call);
\draw[logicline] (call) -- (os);
\draw[logicline] (call) -- (arch);
\draw[logicline] (call) -- (cpu);
\draw[logicline] (os) |- (report);
\draw[logicline] (arch) -- (report);
\draw[logicline] (cpu) |- (report);
\end{tikzpicture}
\caption{A1 system-information logic}
\label{fig:a1-logic}
\end{figure}

\clearpage
\section{Agent A2: Date and Time Agent}

The A2 implementation creates a LangChain agent with \texttt{get\_time}. The
tool returns the host's current \texttt{datetime.now()} value as a formatted
string. It does not currently accept an explicit timezone or location, so its
answer represents the runtime host clock.
The date-and-time processing logic is shown in Figure~\ref{fig:a2-logic}.
\AgentArchitecture{Agent A2}{date, time, or timezone question}{resolve the requested time information}{\texttt{A2\_TOOLS}}{precise date/time response}{a2}
\begin{figure}[!htbp]\centering
\begin{tikzpicture}[node distance=10mm, logiclabel]
\node[logicstep] (request) {Date or time request};
\node[logicstep, right=of request] (invoke) {Invoke \texttt{get\_time}};
\node[logicdata, right=of invoke] (clock) {Runtime host clock\\\texttt{datetime.now()}};
\node[logicstep, right=of clock] (return) {Return formatted date/time string};
\draw[logicline] (request) -- (invoke);
\draw[logicline] (invoke) -- (clock);
\draw[logicline] (clock) -- (return);
\end{tikzpicture}
\caption{A2 date-and-time logic}
\label{fig:a2-logic}
\end{figure}

\clearpage
\section{Agent A3: Research Agent}

The A3 implementation uses two tools: \texttt{wikipedia\_search} calls the
Wikipedia search and summary APIs and returns article titles, extracts, and
URLs; \texttt{trusted\_web\_search} restricts search results to configured
academic and scientific domains, validates result URLs, removes duplicates,
and returns source snippets. The agent prompt requires concise evidence and
prohibits fabricated citations or URLs.
The research and source-verification flow is shown in Figure~\ref{fig:a3-logic}.
\AgentArchitecture{Agent A3}{research question}{search, verify, and extract factual evidence}{Wikipedia and trusted web-search tools}{concise facts with source URL}{a3}
\begin{figure}[!htbp]\centering
\begin{tikzpicture}[node distance=9mm and 14mm, logiclabel]
\node[logicstep] (query) {Research question};
\node[logicstep, below left=of query] (wiki) {Wikipedia API search and article summary};
\node[logicstep, below right=of query] (web) {Trusted-domain web search};
\node[logicdata, below=of wiki] (wresult) {Titles, extracts, and URLs};
\node[logicdata, below=of web] (eresult) {Filtered academic source snippets};
\node[logicstep, below=18mm of query] (verify) {Validate domains, remove duplicate URLs, and select evidence};
\node[logicstep, below=of verify] (answer) {Concise factual answer with sources};
\draw[logicline] (query) -- (wiki);
\draw[logicline] (query) -- (web);
\draw[logicline] (wiki) -- (wresult);
\draw[logicline] (web) -- (eresult);
\draw[logicline] (wresult) |- (verify);
\draw[logicline] (eresult) |- (verify);
\draw[logicline] (verify) -- (answer);
\end{tikzpicture}
\caption{A3 research and source-verification logic}
\label{fig:a3-logic}
\end{figure}

\clearpage
\section{Agent A4: File and Document Agent}

The A4 implementation exposes file discovery, bounded document reading,
TXT/DOCX creation, append and exact replacement, folder creation, rename,
move, metadata, directory listing, document comparison, and heuristic
classification. PDF and DOCX readers are supported through \texttt{pypdf} and
\texttt{python-docx}. There is deliberately no direct delete tool:
\texttt{request\_delete\_confirmation} records a fingerprint in Redis,
publishes a frontend confirmation event, and uses a 300-second expiry before
the backend performs or cancels the operation.
The file-operation and deletion-confirmation flow is shown in Figure~\ref{fig:a4-logic}.
\AgentArchitecture{Agent A4}{file path and requested operation}{inspect or transform the requested file; protect deletion}{file/document tools and confirmation service}{file result or confirmation state}{a4}
\begin{figure}[!htbp]\centering
\begin{tikzpicture}[node distance=9mm and 13mm, logiclabel]
\node[logicstep] (request) {File or document request};
\node[logicdecision, below=of request] (delete) {Deletion requested?};
\node[logicstep, below left=of delete] (normal) {Search, read, create, modify, compare, classify, or manage path};
\node[logicstep, below right=of delete] (inspect) {Resolve path and compute object fingerprint};
\node[logicwarning, below=of inspect] (confirm) {Create Redis confirmation record and emit confirmation event};
\node[logicdata, below=of confirm] (decision) {Wait for backend decision: deleted, cancelled, failed, or expired};
\node[logicdata, below=of normal] (result) {Return file operation result};
\draw[logicline] (request) -- (delete);
\draw[logicline] (delete) -- node[above left] {no} (normal);
\draw[logicline] (delete) -- node[above right] {yes} (inspect);
\draw[logicline] (normal) -- (result);
\draw[logicline] (inspect) -- (confirm);
\draw[logicline] (confirm) -- (decision);
\end{tikzpicture}
\caption{A4 file-operation and deletion-confirmation logic}
\label{fig:a4-logic}
\end{figure}

\clearpage
\section{Content Creator Agent}

The Content Creator implementation creates a tool-free LangChain agent. It
loads the current user name and email for sender context, selects the requested
genre, and drafts, rewrites, summarises, or formats text. It is intentionally
not an email-delivery component and must not send messages or emit unresolved
placeholders.
The content-creation flow is shown in Figure~\ref{fig:content-creator-logic}.
\AgentArchitecture{Content Creator Agent}{writing request and sender context}{choose genre, tone, structure, and format}{language model without external tools}{completed text without placeholders}{content-creator}
\begin{figure}[!htbp]\centering
\begin{tikzpicture}[node distance=9mm and 13mm, logiclabel]
\node[logicstep] (request) {Content request};
\node[logicdecision, below=of request] (type) {Requested genre?};
\node[logicdata, below left=of type] (written) {Email, report, message, article, summary, or documentation};
\node[logicdata, below right=of type] (custom) {Explicit user format and tone};
\node[logicstep, below=18mm of type] (draft) {Draft, rewrite, structure, and remove placeholders};
\node[logicstep, below=of draft] (output) {Return completed text only};
\draw[logicline] (request) -- (type);
\draw[logicline] (type) -- (written);
\draw[logicline] (type) -- (custom);
\draw[logicline] (written) |- (draft);
\draw[logicline] (custom) |- (draft);
\draw[logicline] (draft) -- (output);
\end{tikzpicture}
\caption{Content Creator drafting logic}
\label{fig:content-creator-logic}
\end{figure}

\clearpage
\section{Email Agent}

The Email Agent uses \texttt{send\_email}. The tool reads SMTP host, port,
username, password, TLS, and sender settings from the environment. It sends
through SMTP when credentials are available; otherwise it returns an explicit
simulation-mode response and does not send network mail. The agent extracts
the recipient, subject, and body, removes unresolved placeholders, and returns
the tool's delivery status.
The email-delivery and simulation flow is shown in Figure~\ref{fig:email-logic}.
\AgentArchitecture{Email Agent}{recipient, subject, body, or prior draft}{clean and prepare a deliverable email}{\texttt{EMAIL\_TOOLS} and \texttt{send\_email}}{delivery confirmation}{email}
\begin{figure}[!htbp]\centering
\begin{tikzpicture}[node distance=9mm and 13mm, logiclabel]
\node[logicstep] (request) {Email task and previous draft};
\node[logicstep, below=of request] (fields) {Extract recipient, subject, and body};
\node[logicstep, below=of fields] (clean) {Remove unresolved placeholders};
\node[logicdecision, below=of clean] (smtp) {SMTP credentials available?};
\node[logicdata, below left=of smtp] (send) {Send through SMTP};
\node[logicdata, below right=of smtp] (simulate) {Return safe simulation notice};
\node[logicstep, below=18mm of smtp] (result) {Return delivery status};
\draw[logicline] (request) -- (fields);
\draw[logicline] (fields) -- (clean);
\draw[logicline] (clean) -- (smtp);
\draw[logicline] (smtp) -- node[above left] {yes} (send);
\draw[logicline] (smtp) -- node[above right] {no} (simulate);
\draw[logicline] (send) |- (result);
\draw[logicline] (simulate) |- (result);
\end{tikzpicture}
\caption{Email delivery and simulation-mode logic}
\label{fig:email-logic}
\end{figure}

\clearpage
\section{Weather Agent}

\begingroup
\setlength{\emergencystretch}{3em}
The Weather Agent uses Open-Meteo through \texttt{get\_current\_weather} and
\texttt{get\_weather\_forecast}. Each operation geocodes the city first, then
requests current or daily forecast data, maps WMO weather codes to readable
conditions, and reports Celsius/Fahrenheit temperatures, humidity, wind, and
precipitation. Forecast length is bounded to one through seven days.
\par
\endgroup
The weather-retrieval and reporting logic is shown in Figure~\ref{fig:weather-logic}.
\AgentArchitecture{Weather Agent}{location and weather request}{resolve location and requested time horizon}{current-weather and forecast tools}{weather report in clear units}{weather}
\begin{figure}[!htbp]\centering
\begin{tikzpicture}[node distance=9mm and 13mm, logiclabel]
\node[logicstep] (request) {Weather request and city};
\node[logicstep, below=of request] (geo) {Geocode city with Open-Meteo};
\node[logicdecision, below=of geo] (kind) {Current conditions or forecast?};
\node[logicdata, below left=of kind] (current) {Current temperature, feels-like, humidity, wind, precipitation};
\node[logicdata, below right=of kind] (forecast) {Daily high, low, weather code, and precipitation};
\node[logicstep, below=18mm of kind] (format) {Map WMO codes and format Celsius/Fahrenheit report};
\draw[logicline] (request) -- (geo);
\draw[logicline] (geo) -- (kind);
\draw[logicline] (kind) -- node[above left] {current} (current);
\draw[logicline] (kind) -- node[above right] {forecast} (forecast);
\draw[logicline] (current) |- (format);
\draw[logicline] (forecast) |- (format);
\end{tikzpicture}
\caption{Weather retrieval and reporting logic}
\label{fig:weather-logic}
\end{figure}

\clearpage
\section{Loan EMI Calculator Agent}

The Loan EMI Agent uses \texttt{compute\_emi\_and\_schedule}. It converts the
annual percentage rate to a monthly rate, computes the fixed-payment EMI,
derives total payment and interest, and produces a monthly schedule for short
tenures or an annual amortisation breakdown for longer tenures.
The loan-EMI and amortisation logic is shown in Figure~\ref{fig:loan-emi-logic}.
\AgentArchitecture{Loan EMI Calculator Agent}{principal, rate, and tenure}{validate financial parameters and calculate repayment}{\texttt{LOAN\_EMI\_TOOLS}}{EMI and amortisation result}{loan-emi}
\begin{figure}[!htbp]\centering
\begin{tikzpicture}[node distance=9mm and 13mm, logiclabel]
\node[logicstep] (input) {Principal, annual rate, and tenure};
\node[logicdata, below=of input] (emi) {Convert annual rate to monthly rate and compute fixed EMI};
\node[logicstep, below=of emi] (totals) {Compute total payment and total interest};
\node[logicdecision, below=of totals] (term) {Tenure at most 12 months?};
\node[logicdata, below left=of term] (monthly) {Monthly amortisation table};
\node[logicdata, below right=of term] (annual) {Annual amortisation breakdown};
\node[logicstep, below=18mm of term] (result) {Return calculation summary and schedule};
\draw[logicline] (input) -- (emi);
\draw[logicline] (emi) -- (totals);
\draw[logicline] (totals) -- (term);
\draw[logicline] (term) -- node[above left] {yes} (monthly);
\draw[logicline] (term) -- node[above right] {no} (annual);
\draw[logicline] (monthly) |- (result);
\draw[logicline] (annual) |- (result);
\end{tikzpicture}
\caption{Loan EMI and amortisation logic}
\label{fig:loan-emi-logic}
\end{figure}

\clearpage
\section{Crypto PnL Agent}

The Crypto PnL Agent uses \texttt{calculate\_crypto\_pnl}. It calculates buy
and sell notional values, applies the fee percentage to both sides, derives
net profit, and computes ROI relative to total purchase outflow.
The cryptocurrency PnL and ROI logic is shown in Figure~\ref{fig:crypto-pnl-logic}.
\AgentArchitecture{Crypto PnL Agent}{trade details and report request}{calculate net PnL and return on investment}{\texttt{CRYPTO\_PNL\_TOOLS}}{formatted investment report}{crypto-pnl}
\begin{figure}[!htbp]\centering
\begin{tikzpicture}[node distance=9mm and 13mm, logiclabel]
\node[logicstep] (input) {Quantity, buy price, sell price, and fee rate};
\node[logicdata, below left=of input] (buy) {Buy notional plus purchase fee};
\node[logicdata, below right=of input] (sell) {Sell notional minus sale fee};
\node[logicstep, below=18mm of input] (profit) {Net profit = net proceeds minus total outflow};
\node[logicstep, below=of profit] (roi) {ROI = net profit divided by total outflow};
\node[logicstep, below=of roi] (report) {Return PnL and ROI report};
\draw[logicline] (input) -- (buy);
\draw[logicline] (input) -- (sell);
\draw[logicline] (buy) |- (profit);
\draw[logicline] (sell) |- (profit);
\draw[logicline] (profit) -- (roi);
\draw[logicline] (roi) -- (report);
\end{tikzpicture}
\caption{Cryptocurrency PnL and ROI logic}
\label{fig:crypto-pnl-logic}
\end{figure}

\clearpage
\section{Code Runner Agent}

The Code Runner Agent uses \texttt{build\_run\_python\_code\_tool}. It limits
source code to 20,000 characters, executes Python with isolated mode and a
sanitised environment, bounds runtime to ten seconds, captures at most 64,000
output characters, publishes terminal events, and limits concurrent runs to
two. The agent repairs failures for at most three attempts.
The code-generation, sandbox-execution, and repair logic is shown in Figure~\ref{fig:code-runner-logic}.
\AgentArchitecture{Code Runner Agent}{programming request and prior feedback}{generate, execute, inspect, and repair Python}{sandboxed Python execution tool}{stdout, duration, or bounded failure report}{code-runner}
\begin{figure}[!htbp]\centering
\resizebox{0.82\textwidth}{!}{\begin{tikzpicture}[node distance=9mm and 13mm, logiclabel]
\node[logicstep] (request) {Coding or testing request};
\node[logicstep, below=of request] (generate) {Generate standalone Python script};
\node[logicstep, below=of generate] (execute) {Run isolated Python subprocess};
\node[logicdecision, below=of execute] (success) {Exit code zero?};
\node[logicdata, below left=of success] (done) {Return stdout and duration};
\node[logicdecision, below right=of success] (attempts) {Fewer than three attempts?};
\node[logicstep, below=of attempts, xshift=-20mm] (repair) {Repair script using traceback};
\node[logicwarning, below right=of attempts] (failed) {Return failure report};
\draw[logicline] (request) -- (generate);
\draw[logicline] (generate) -- (execute);
\draw[logicline] (execute) -- (success);
\draw[logicline] (success) -- node[above left] {yes} (done);
\draw[logicline] (success) -- node[above right] {no} (attempts);
\draw[logicline] (attempts) -- node[right] {yes} (repair);
\draw[logicline] (repair.west) -- ++(-60mm,0) |- ([xshift=-3mm]generate.west);
\draw[logicline] (attempts) -- node[above] {no} (failed);
\end{tikzpicture}}
\caption{Code generation, sandbox execution, and repair logic}
\label{fig:code-runner-logic}
\end{figure}

\clearpage
\section{Agent Creator}

The Agent Creator implementation is responsible for unsupported capabilities.
\texttt{save\_tool\_file} parses generated source with the Python AST before
writing and registering it. \texttt{create\_and\_register\_agent} generates
the agent module, updates the agent registry and planner constants, validates
the changes, and rolls them back on failure. \texttt{run\_generated\_agent}
then loads and executes the new agent. The implementation is present in the
source directory but is commented out of the default registry.
The dynamic agent-creation and registration logic is shown in Figure~\ref{fig:agent-creator-logic}.
\AgentArchitecture{Agent Creator}{unsupported specialised task}{design capability, create files, and register new components}{tool-file, agent-file, and registry tools}{result from the generated agent}{agent-creator}
\begin{figure}[!htbp]\centering
\begin{tikzpicture}[node distance=9mm and 13mm, logiclabel]
\node[logicstep] (request) {Unsupported specialised capability};
\node[logicstep, below=of request] (tool) {Generate tool source, AST-validate, save, and register};
\node[logicstep, below=of tool] (agent) {Generate agent source and update registries};
\node[logicdecision, below=of agent] (valid) {Registration succeeds?};
\node[logicwarning, below left=of valid] (rollback) {Restore registry and agent files};
\node[logicstep, below right=of valid] (run) {Load and execute generated agent};
\node[logicdata, below=18mm of valid] (result) {Return generated-agent result};
\draw[logicline] (request) -- (tool);
\draw[logicline] (tool) -- (agent);
\draw[logicline] (agent) -- (valid);
\draw[logicline] (valid) -- node[above left] {no} (rollback);
\draw[logicline] (valid) -- node[above right] {yes} (run);
\draw[logicline] (run) |- (result);
\end{tikzpicture}
\caption{Dynamic agent creation and registration logic}
\label{fig:agent-creator-logic}
\end{figure}

